\scanpage{213}

at least \inserted{we will have} a partial splitting over
$\Gamma_{\mathbb Q}(q)$, which we will try to determine ;
$\mathcal N_{0,3}^{!\sim}(q)$ denotes a hypothetical
subgroup of $\mathcal N_{0,3}^{!\sim}$
(which we would have to \unsure{manage} to determine, \unsure{if indeed}
\illegible{1} not too \illegible{1}...), inducing on $\Gamma_{\mathbb Q}$
the \inserted{open} subgroup $\Gamma_{\mathbb Q}(q)$ of
$\Gamma_{\mathbb Q}$ stabilizer of $q$.

\margin{The definition is forced for example when $q$ is
an \unsure{element} of a cyclotomic extension of $\mathbb Q$, in such
a way that $\mathbb Q(q)$ corresponds to a finite discrete
quotient group of $\widehat{\mathbb Z}^{*}$ and $\Gamma_{\mathbb Q}(q)$ is
the kernel of
\diagram{\begin{tikzcd}[ampersand replacement=\&]
\Gamma_{\mathbb Q}\arrow[r,"\chi"]\&\widehat{\mathbb Z}^{*}\arrow[r]\&G
\end{tikzcd}}\ldots}

The analysis of cases $1^\circ$ and $2^\circ$ above in fact amounted to
the determination of such splittings, for $q=\frac12$ and for
$q=-\overline j$ --- in the first case $\Gamma_{\mathbb Q}(q)$ is
the whole of $\Gamma_{\mathbb Q}$, in the second it is the subgroup
of index $2$ kernel of $\varepsilon_2\circ\chi:\Gamma_{\mathbb Q}
\longrightarrow\{\pm1\}$, and idem for
$\unsure{\mathcal N|\mathcal M}_{0,3}^{!\sim}(q)$.
By transforming by $\rho$ \inserted{and $\rho^2$}, we
\deleted{\unsure{would lead}}\inserted{would deduce} \unsure{from the splittings|from the splitting} $k_{1/2}$
the splittings $k_2$ and $k_{-1}$ --- by transforming by any one
of the $\sigma_{\unsure{i|\rho}}$,
\deleted{or by one of the $\unsure{\tau_i|\sigma_i}$}, we would deduce
from the splitting $k_{-\overline j}$ a splitting of type $k_{\unsure{-j|-1}}$.

As for $\mathcal M_{0,3}^{!\sim}[0]$, it corresponds
\inserted{in principle} to the point

\scanpage{214}
\mspage{406}

$0$ of $\widehat{\mathcal U}_{0,3}=\unsure{\mathbb P^1_{\mathbb Q}}$,
there \deleted{\unsure{correspond to it by}}\inserted{two other}
analogous subgroups $\mathcal M_{0,3}^{!\sim}[q]$, with
$q=1$, $q=\infty$, corresponding to the two other points ``at infinity''
of $\mathcal U_{0,3}$, and which are deduced from the preceding one by applying
$\rho$, $\rho^2$.
Here, we have groups which \uline{do} \uline{not} correspond
\inserted{\unsure{strictly}} to \inserted{even} partial sections
of $\mathcal M_{0,3}^{!\sim}$ over $\mathcal N_{0,3}^{!\sim}$,
since the intersection with $\Pi_{0,3}^{\wedge}$ is not reduced to $1$,
but it is one of the three loop groups of type $L_0,L_1,L_\infty$.

But we know that \inserted{the choice of }a \unsure{uniformizer|uniformizer}
of the local ring of $\widehat{\mathcal U}_{0,3}$ at $0$ (say)
determines a splitting of
$\mathcal M_{0,3}^{!\unsure{\sim}}\unsure{(0)|[0]}$, regarded as
an extension of $\mathcal N_{0,3}^{!}$ by $L_0$ --- and later on we will have to
try to grasp \deleted{\unsure{this}}\inserted{such a}
splitting \inserted{(associated with a \unsure{uniformizer})}
from the viewpoint of our calculation, as a \unsure{privileged} way
\deleted{of ``\unsure{normalizing}''}, for \unsure{every}
\deleted{$\alpha,\beta$} \unsure{datum} \inserted{mod $L_0$} of $(\alpha,\beta)$
satisfying a loop relation, of ``normalizing'' the choice of the
$(\alpha,\beta)$ mod $L_0$.
I will come back to this later, surely.

\scanpage{215}
\mspage{407}

\uline{V} \uline{Recapitulation} \uline{in terms of}
$\mathcal M_{0,3}^{!\sim}[0]$

It finally appears that the simplest way to study the loop
automorphisms of $\Pi_{0,3}^{\wedge}$ (commuting outer\textsuperscript{ly}
with $\mathfrak S_3$, i.e. the elements of \deleted{\illegible{1}}
$\mathcal N_{0,3}^{!\unsure{\sim}}$) is to begin with the study of the
subgroup $\mathcal M_{0,3}^{!\unsure{\sim}}[0]$, normalizer
of $L_0$ \inserted{in $\mathcal M_{0,3}^{!\sim}$}.
Setting for $u\in\mathcal M_{0,3}^{!\unsure{\sim}}[0]$
\[
\tag{128}
u(\sigma_\infty)=\eta\sigma_\infty,
\qquad u(\rho)=\xi\rho
\qquad(\eta,\xi\in\Pi_{0,3}^{\wedge})
\]
i.e.
\[
\tag{129}
[u,\sigma_\infty]=\eta\qquad [u,\rho]=\xi
\]
which express the outer commutation of $u$ with $\mathfrak S_3$
(indeed generated by $\sigma_\infty$ and $\rho$), the condition of
normalizing $L_0$ is expressed by a relation of the form
\[
\tag{130}
u(\varepsilon_0)=\varepsilon_0^\mu
\quad\text{or}\quad \unsure{u}(\ell_0)=\ell_0^\mu
\qquad(\mu\in\widehat{\mathbb Z}^{*}).
\]

\margin{(130) is also equivalent, given $[u,\rho]=\xi$, to
\[
\tag{130bis}
\left\{
\begin{aligned}
u(\varepsilon_1)&=\operatorname{int}(\xi)(\varepsilon_1^\mu)\\
&\text{and to}\\
u(\varepsilon_\infty)&=
\operatorname{int}\bigl(\rho^{-1}(\xi^{-1})\,\illegible{1}\bigr)(\varepsilon_\infty^\mu)
\end{aligned}
\right.
\]}

\inserted{Now we have, since $\varepsilon_0=\sigma_\infty\rho$}
\deleted{which, taking $\varepsilon_0=\sigma_\infty\rho$ into account,
takes (by an immediate calculation) the form}
\deleted{\unsure{is written}}
\deleted{(131)} \deleted{$\unsure{\eta\sigma_\infty}$}
\[
u(\varepsilon_0)=u(\sigma_\infty)u(\rho)
=\eta\sigma_\infty\xi\rho
=\eta\sigma_\infty(\xi)\underbrace{\sigma_\infty\rho}_{\varepsilon_0}
=\eta\sigma_\infty(\xi)\varepsilon_0
\]
so (130) can be written
\[
\eta\sigma_\infty(\xi)=\varepsilon_0^{\mu-1}
\qquad\text{i.e., setting }\mu-1=2\nu
\quad(\nu\in\widehat{\mathbb Z})
\]
\[
\eta\sigma_\infty(\xi)=\ell_0^\nu
\qquad\text{i.e., applying }\sigma_\infty
\]
\[
\tag{131}
\sigma_\infty(\eta)\xi=\ell_1^\nu.
\]
Now the relations
\[
\tag{132}
(\eta\sigma_\infty)^2=1,\qquad(\xi\rho)^3=1
\]
can also be written
\[
\tag{133}
\boxed{\eta\sigma_\infty(\eta)=1,\qquad
\xi\rho(\xi)\rho^2(\xi)=1},
\]
whence $\sigma_\infty(\eta)=\eta^{-1}$, so that (131) can
also be written in a more \uline{convenient} way
\[
\tag{134}
\boxed{\xi=\eta\ell_1^\nu}
\qquad(\text{loop relation})
\]

\scanpage{216}

We thus find the

\uline{Theorem} The elements $u$ of
$\unsure{\mathcal M|\mathcal N}_{0,3}^{!\sim}[0]$ --- i.e. the
outer automorphisms of $\Pi_{0,3}^{\wedge}$ which commute outer\textsuperscript{ly}
with $\mathfrak S_3$ and which normalize $L_0$ --- are in one-to-one
correspondence \inserted{(via formulas (128) or 129)}
with the pairs $(\xi,\eta)$ $(\xi,\eta\in\Pi_{0,3}^{\wedge})$,
satisfying relations (133)(134), and such that
\inserted{moreover (``preliminary condition'')}
relations (132) (equivalent : (133)) are
\deleted{\unsure{primitive}}\inserted{generating}
for $\widehat{\mathcal G}_{0,3}^{+}$.
(\deleted{\illegible{1}} In (134) $\nu$ is a suitable element of
$\widehat{\mathbb Z}$, obviously uniquely determined, and subject
to the additional condition $2\nu+1\in\widehat{\mathbb Z}^{*}$)
(which is implied by the preliminary condition, and perhaps implies it...)

\margin{this condition $2\nu+1\in\widehat{\mathbb Z}^{*}$ is in fact
a consequence of equations (133), (134), cf. below, theorem page 409!}

\unsure{And} the obvious inclusion
\[
\tag{135}
\begin{aligned}
L_0&\hookrightarrow\mathcal M_{0,3}^{!\unsure{\sim}}\unsure{(0)|(\omega)}\\
\ell_0^n=\lambda&\longmapsto\operatorname{int}(\lambda)
\end{aligned}
\]
is given by
\[
\tag{136}
\ell_0^n=\lambda\longmapsto
\bigl(\xi=\lambda\rho(\lambda)^{-1}=\ell_0^n\ell_1^{-n},
\quad\eta=\lambda\sigma_\infty(\lambda)^{-1}=\ell_0^n\ell_1^{-n}\bigr)
\]
and more generally the left action of $L_0$ on
$\mathcal M_{0,3}^{!\unsure{\sim}}\unsure{(0)|(\omega)}$ is given,
for $\lambda=\ell_0^n\in L_0$, by
\[
\tag{137}
(\xi,\eta)\longmapsto
\bigl(\xi'=\ell_0^n\xi\ell_1^{-n},\quad
\eta'=\ell_0^n\eta\ell_1^{-n}\bigr)
\]
so
\[
\mathcal M_{0,3}^{!\unsure{\sim}}[0]/L_0
\simeq\mathcal N_{0,3}^{!\unsure{\sim}}
\]
is isomorphic to the quotient of the set of pairs $(\xi,\eta)$
satisfying (133)(134) [and the preliminary condition], by the operations
(137) --- taking into account a \deleted{\unsure{operation}}
natural group structure : to put on this set of pairs
$(\xi,\eta)$, given by

\scanpage{217}
\mspage{408}

\[
\tag{138}
\underbrace{(\xi',\eta')}_{\deleted{\unsure{u'}}}(\xi,\eta)=(\xi'',\eta'')
\]
with
\[
\tag{139}
\xi''=u'(\xi)\xi',\qquad \eta''=u'(\eta)\eta'
\]
$u'$ being the automorphism defined by $(\xi',\eta')$.

\uline{Corollary} Every $u\in\unsure{\mathcal M_{0,3}^{!\sim}}$ can be written \unsure{under|in} the form
\[
\tag{140}
u=f^{-1}u_{\xi,\eta}=\operatorname{int}(f^{-1})\circ u_{\xi,\eta},
\]
with $f\in\widehat{\Pi}_{0,3}$ and $(\xi,\eta)$ satisfying the conditions of the preceding theorem, so that we will have
\[
\tag{141}
u(\unsure{\sigma_\infty})=h\sigma_\infty,\qquad u(\rho)=g\rho,\quad\text{with}
\]
\[
\tag{143}
h=f^{-1}\eta\sigma_\infty(f),\qquad g=f^{-1}\xi\rho(f).
\]
For us to have
\[
f^{-1}u_{\xi,\eta}=f'^{-1}u_{\xi',\eta'}
\]
it \unsure{is necessary and sufficient} that there exist $n\in\mathbb Z$ such that
\[
\tag{144}
\xi'=\ell_0^n\xi\ell_1^{-n},\qquad
\eta'=\ell_0^n\eta\ell_1^{-n},\qquad f'=\ell_0^n f.
\]

\uline{N.B.} \unsure{Thanks} to the loop formula, once the multiplier $\mu$ i.e. $\nu$ is fixed (and in fact $\mu$ is already uniquely determined by the knowledge of $\xi$ alone, as is \unsure{shown} by the adelic calculations in $\mathrm{Gl}(2,\widehat{\mathbb Z})'\ldots$), $\xi$ and $\eta$ determine one another \inserted{by (134)}, so we can eliminate one or the other from the system of equations (133), (134), which becomes either
\[
\tag{145}
\begin{gathered}
\xi\rho(\xi)\rho^2(\xi)=1,\qquad
(\deleted{\unsure{\eta}}\inserted{\xi}\ell_1^{-\nu})\sigma_\infty(\xi\ell_1^{-\nu})=1
\quad\text{i.e.}\\
\sigma_\infty(\xi)=\ell_1^\nu\xi^{-1}\ell_0^\nu
\end{gathered}
\]
or
\[
\tag{146}
\eta\sigma_\infty(\eta)=1,\qquad
(\eta\ell_1^\nu)\rho(\eta\ell_1^\nu)\rho^2(\eta\ell_1^{\unsure{\nu|-\nu}})=1.
\]

\scanpage{218}

The adelic calculation. We again denote by $\xi,\eta$ the images of $\xi,\eta$ in $\mathrm{Sl}(2,\widehat{\mathbb Z})'$, which lift to elements, again denoted $\xi,\eta$, of $\mathrm{Sl}(2,\widehat{\mathbb Z})$.

There is a sign ambiguity, which we remove for $\xi$ by imposing
\[
\tag{147}
\xi\rho(\xi)\rho^2(\xi)=\uuline{+}1
\qquad(\text{in }\mathrm{Sl}(2,\widehat{\mathbb Z}))
\]
\inserted{for $\eta$} there \uline{is no simple way} to remove the \unsure{ambiguity} (if we wish to admit \inserted{provisionally} the possibility of writing in the form $\beta\sigma_\infty(\beta)^{-1}$) \uline{\unsure{of which|thus}} we will have a priori only
\[
\tag{148}
\eta\sigma_\infty(\eta)=\deleted{\unsure{\varepsilon}}\pm1,
\]
with no possibility of changing the sign in the $2^\circ$ case by the change $\eta\mapsto-\eta$ ! Let us set

\margin{\uline{but yes}, we can make the choice of $\eta$ by requiring $\xi=+\eta\ell_1^\nu$ (but this convention \unsure{is} not compatible with a convention \unsure{made} elsewhere $\eta=\beta\sigma_\infty(\beta)^{-1}$).}
\margin{$\in\widehat{\Pi}_{0,3}$}

\[
\tag{148}
\xi=\begin{pmatrix}A&B\\C&D\end{pmatrix},\qquad
\eta=\begin{pmatrix}R&S\\T&U\end{pmatrix}.
\]
We will have
\[
\tag{149}
AD-BC=1\qquad RU-ST=1
\]
and relation (147) is then equivalent to :
\[
\tag{150}
C+\unsure{D}-B=1,
\]
\deleted{\unsure{whereas}} (at least in each $\ell$-adic component we will have to have either this relation, or
\[
\xi_\ell=\begin{pmatrix}A_\ell&B_\ell\\C_\ell&D_\ell\end{pmatrix}
=-\rho^{-1}=\begin{pmatrix}1&-1\\-1&0\end{pmatrix},
\]
\deleted{\unsure{cf.}} proposition on \unsure{page} 394 -- but this solution is incompatible with the condition $\boxed{\xi\equiv1\ (2)}$, whence $A,D\equiv1\ (2)$ and $B,C\equiv0\ (2)\ldots$ so (150) is true in $\mathbb Z_\ell$ for $\forall\ell$, hence in $\widehat{\mathbb Z}\ldots$).

\scanpage{219}
\mspage{409}

Relation (148) can also be written
\[
\eta=\pm\sigma_\infty(\eta^{-1})
\]
now $\sigma_\infty(\eta^{-1})={}^t\eta$, so (148) can be written
\[
\eta=\pm{}^t\eta,
\]
but if we had the sign $-$, we would have $U=-U$ so $2U=0$ so $U=0$, which is incompatible with
\[
\eta\equiv\unsure{\pm1|1}\quad(2)
\]
which implies $R,U\equiv1\ (2)$, $S,T\equiv0\ (2)$. We will therefore have $\eta={}^t\eta$ i.e.
\[
\tag{151}
\eta\sigma_\infty(\eta)=+1,\qquad\text{i.e.}
\]
\[
\tag{152}
S=T.
\]

The loop relation (134) can be written
\[
\tag{153}
\begin{aligned}
\begin{pmatrix}A&B\\C&D\end{pmatrix}
&=\deleted{\begin{pmatrix}1&0\\2\nu&1\end{pmatrix}}
\inserted{\varepsilon}\begin{pmatrix}R&S\\T&U\end{pmatrix}
\begin{pmatrix}1&0\\2\nu&1\end{pmatrix}\\
&=\varepsilon\begin{pmatrix}R+2\nu S&S\\T+2\nu U&U\end{pmatrix}
\qquad\text{with }\varepsilon\in\{\pm1\}
\end{aligned}
\]

\margin{\unsure{Observation} here $\varepsilon=1$, by choosing the sign of $\eta$ suitably (153) (the sign convention does not coincide with a previous convention centered on the $\alpha,\beta$, but from the present viewpoint it seems more \unsure{friendly}).}

and it is equivalent to the relations
\[
\tag{154}
B=\varepsilon S,\qquad D=\varepsilon U\qquad\text{and}
\]
\deleted{\illegible{1}}
\[
\tag{155}
\mu D=1\qquad\text{where }\mu=2\nu+1
\]
\unsure{equivalent|being equivalent} also ; given (154), to
\[
\tag{155bis}
\mu U=\varepsilon,
\]
and the latter merely expresses \inserted{(taking (152) into account)} the relation
\[
\varepsilon\overbrace{(T+2\nu U)}^{A'}
+\overbrace{\varepsilon U}^{B'}-\overbrace{\varepsilon S}^{C'}=1
\]
\unsure{paraphrasing}
\[
A+B-C=1,
\]
it